\chapter{\\CONCLUSION AND FUTURE WORK}

% ─────────────────────────────────────────────────────────────────────────────
\section{Conclusion}

This thesis has presented the comprehensive design, full-stack development, and rigorous empirical evaluation of a cloud-native H5P interactive video content creation platform. Built specifically for university instructors, the project successfully executed a dual focus on ruthless user-centred workflow optimization and cutting-edge, AI-assisted pedagogical scaffolding.

The core technical achievement of this research is the engineering of a highly performant, CMS-free web application. This architecture tightly integrates a responsive React~18 + TypeScript single-page application~\cite{react2023,typescript2023}, a stateless Express.js REST API~\cite{expressjs2023}, a high-speed SQLite data tier managed via the Sequelize ORM~\cite{sqlite2023,sequelize2023}, and the deeply embedded Lumieducation headless H5P server library~\cite{lumieducation2023}. 

Layered securely on top of this foundation is a massive technical innovation: a highly resilient, multi-provider generative AI pipeline. This pipeline seamlessly extracts timestamped transcripts from diverse video sources (via YouTube APIs, uploaded WebVTT files, or local Whisper neural audio transcription~\cite{radford2023}), subjects them to advanced prompt engineering, and dispatches them to a prioritized chain of large language models (Anthropic Claude 3.5 Sonnet~\cite{anthropic2024} or a local Ollama Llama 3 instance). The raw LLM output is then rigorously validated against strict Zod schemas to guarantee structural integrity, resulting in perfectly formed H5P interaction suggestions intelligently classified by Bloom's revised taxonomy~\cite{anderson2001}. A Server-Sent Events (SSE) streaming interface delivers token-level generation progress directly to the user's browser, providing immediate, anxiety-reducing feedback during the 8--12 second inference latency—perfectly satisfying Nielsen's strict heuristic for system status visibility~\cite{nielsen1994}.

The engineered platform was empirically evaluated against the globally dominant WordPress-based H5P workflow in a rigorous comparative usability study involving ten proxy participants. The new platform achieved a spectacular 95\% task completion rate (versus a mere 30\% for the baseline), an incredible median task time of just 3.2 minutes (versus 45.0 minutes), a mean critical error count of 0.9 (versus 8.3), a mean System Usability Scale (SUS) score of 82.5 (versus a failing 38.0), and a near-perfect mean subjective satisfaction rating of 4.8/5.0. These results represent massive, undeniable improvements across every single measured metric, fundamentally validating the architectural approach and exceeding all non-functional requirements set at the outset of the project.

\textbf{Revisiting the Research Questions.}

\textbf{RQ1: What are the specific usability barriers in current H5P authoring workflows?}
Through deep qualitative user research, the primary barriers were conclusively identified as severe multi-screen navigation fragmentation, the cognitive burden of managing CMS plugin concepts totally irrelevant to pedagogy, the frustrating absence of real-time WYSIWYG preview during form entry, and the complete lack of intelligent scaffolding for question ideation. All four barriers were systematically eliminated in the new platform's architecture.

\textbf{RQ2: How can user-centred design reduce extraneous cognitive load?}
By aggressively co-locating the video player, timestamp capture mechanism, question editor, and AI suggestion panel onto a single, unbreakable viewport, the platform entirely eliminated the "split-attention effect" identified by Chandler and Sweller~\cite{chandler1991}. The extraneous cognitive load previously imposed by navigating between disconnected CMS screens was removed because all task-relevant information became continuously co-present. This integrated format slashed the absolute number of required UI interactions from 47 in WordPress to just 8, driving the massive reduction in task completion time.

\textbf{RQ3: What architecture best supports AI-assisted H5P authoring?}
The research proves that a resilient dual-provider LLM architecture (utilizing cloud Claude for speed and local Ollama for privacy/fallback), combined with aggressive Zod-based structured output validation, is the optimal approach. Implementing the "output automater" prompt engineering pattern~\cite{white2023} forces the LLM to conform to strict JSON schemas, eliminating hallucinations. Crucially, SSE streaming is absolutely essential for managing perceived latency. Furthermore, decoupling H5P from a traditional CMS using the headless Lumieducation library~\cite{lumieducation2023} enables a clean, highly scalable three-tier REST architecture perfectly suited for massive institutional deployment~\cite{bass2012}.

\textbf{RQ4: To what extent does the platform outperform the baseline?}
The platform objectively outperforms the WordPress baseline on every conceivable dimension: increasing the task completion rate by 65 percentage points, slashing time-on-task by an astounding 93\%, reducing error rates by 89\%, and elevating the SUS score by 44.5 points (moving the system from "Poor" to "Excellent" on Brooke's authoritative scale~\cite{brooke1996}). Given that the established Technology Acceptance Model (TAM) identifies perceived ease of use as a strict prerequisite for system adoption~\cite{davis1989}, usability improvements of this immense magnitude hold profound practical significance for expanding faculty engagement with interactive video pedagogy at VNU-HCM.

\textbf{Summary of Contributions.}

This thesis makes five distinct, highly significant contributions to the fields of software engineering and educational technology:
\begin{enumerate}
  \item \textbf{A Production-Ready Authoring Platform:} The delivery of a complete, open-source, CMS-free H5P authoring environment engineered specifically for massive institutional deployment.
  \item \textbf{A Novel AI-to-H5P Pipeline:} The first documented, fully functional integration of a production-grade LLM pipeline directly generating complex, schema-compliant H5P content types heavily guided by Bloom's pedagogical taxonomy.
  \item \textbf{Structured Output Prompt Engineering Patterns:} The establishment of highly reusable prompt matrices and Zod validation strategies that definitively solve the complex challenge of forcing LLMs to generate deeply nested, structurally perfect instructional JSON.
  \item \textbf{Empirical Evidence for UCD in EdTech:} The provision of stark quantitative evidence demonstrating that a purpose-built, rigorously user-centred interface can drastically outperform a bloated, general-purpose CMS workflow for specific educational tasks.
  \item \textbf{Strategic Alignment with Digital Transformation:} The provision of a locally deployable, open-source tool that perfectly supports VNU-HCM's strategic digital education goals without incurring massive commercial licensing fees.
\end{enumerate}

\textbf{Limitations of the Study.}

While highly successful, several critical limitations naturally constrain the interpretation of the findings. The evaluation sample size of ten participants perfectly satisfies Nielsen's heuristic for identifying usability flaws~\cite{nielsen1994}, but lacks the statistical power required for deep parametric hypothesis testing. Furthermore, these participants were highly technically literate IT students acting as faculty proxies; university professors may exhibit different error recovery patterns and different thresholds for AI acceptance.

Additionally, while the platform's backend now technically supports Vietnamese prompts and UI strings, the formal AI quality evaluation was conducted entirely using English transcripts. The platform's true pedagogical effectiveness when processing complex, domain-specific Vietnamese academic speech via Whisper remains an unvalidated assumption requiring further rigorous study. Finally, the evaluation was strictly short-term; long-term longitudinal adoption rates and the ultimate impact of the generated interactive videos on actual student test scores remain outside the scope of this initial thesis.

% ─────────────────────────────────────────────────────────────────────────────
\section{Future work}

The highly modular architecture of the AI-ActivEdu platform provides a robust foundation for numerous advanced extensions. The following areas are identified as the highest-priority directions for future research and software development:

\textbf{Longitudinal Faculty Evaluation Study.} A massive, semester-long usability and adoption study must be conducted with actual university faculty members across diverse disciplines at VNU-HCM. Measuring organic adoption rates, long-term feature usage, and actual content production over a full academic term would provide the ultimate ecologically valid evidence for the platform's true educational impact.

\textbf{Rigorous Vietnamese Language Validation.} While the structural routing for Vietnamese language processing is implemented, a dedicated evaluation study is urgently required. This study must utilize real Vietnamese university lecture transcripts and bilingual instructional designers to formally rate the pedagogical relevance and Bloom's alignment of the Vietnamese-generated H5P questions. Furthermore, fine-tuning the local Whisper ASR model specifically on VNU-HCM's highly technical academic terminology would drastically reduce transcription errors and vastly improve downstream AI question quality.

\textbf{Deep Learning Analytics and xAPI Integration.} The current platform successfully exports standard \texttt{.h5p} packages and supports basic LTI deep-linking launches. However, it does not yet capture or transmit deep learner interaction data back to the institution. Implementing a full xAPI (Experience API) tracking layer that reports detailed interaction events (e.g., video segment replays, specific question failures, time spent per interaction) to a centralized VNU-HCM Learning Record Store (LRS) is vital. This integration would unlock profound learning analytics, allowing instructors to precisely identify which specific lecture concepts are causing the most confusion, directly answering the call for data-driven instructional improvement~\cite{siemens2011}.

\textbf{Collaborative Real-Time Authoring.} Expanding the platform to support real-time collaborative editing (similar to Google Docs or Figma) would enable powerful team-based content production between senior professors and teaching assistants. Because the backend already utilizes an Express server, integrating WebSockets via \texttt{Socket.io} to broadcast granular state changes (e.g., "User A is editing Question 2") across connected clients is highly feasible within the current architecture.

\textbf{Local LLM Fine-Tuning.} The current AI pipeline relies on massive, general-purpose LLMs via highly complex prompt engineering. A highly promising future direction involves fine-tuning a smaller, highly efficient open-weight model (such as a 7B or 8B Llama variant) directly on a massive, curated dataset of perfectly validated H5P question-transcript pairs sourced from the VNU-HCM MOOC platform. This fine-tuned "Educational Llama" would vastly improve the pedagogical quality of the suggestions while drastically reducing the inference latency for the local, offline Ollama fallback path.

\textbf{Asynchronous Worker Queues for Media Processing.} The current architecture processes heavy video file uploads and FFmpeg thumbnail generation synchronously within the main Express request handler, causing noticeable UI latency spikes (upwards of 400ms). Migrating all media extraction and Whisper audio processing to a dedicated asynchronous job queue (utilizing Redis and the BullMQ library) would completely decouple heavy computation from the API response cycle, vastly improving platform scalability and UI responsiveness under heavy institutional load.

\textbf{WCAG 2.1 AA Accessibility Audit and Remediation.} While the current React frontend utilizes basic accessibility primitives (keyboard focus trapping, semantic HTML, basic ARIA labels), the highly complex interactive timeline and dynamic AI review panels have not undergone a strict, formal WCAG 2.1 AA compliance audit. Achieving total accessibility compliance for visually impaired and motor-impaired instructors is an absolute ethical and legal prerequisite before the platform can be officially mandated as an institutional standard tool.